\documentclass[conference]{IEEEtran}
\IEEEoverridecommandlockouts

\usepackage{cite}
\usepackage{amsmath,amssymb,amsfonts}
\usepackage{algorithmic}
\usepackage{graphicx}
\usepackage{textcomp}
\usepackage{xcolor}
\usepackage{booktabs}
\usepackage{array}
\usepackage{multirow}

\def\BibTeX{{\rm B\kern-.05em{\sc i\kern-.025em b}\kern-.08em
    T\kern-.1667em\lower.7ex\hbox{E}\kern-.125emX}}

\begin{document}

% ─────────────────────────────────────────────────────────────────────────────
%  TITLE
% ─────────────────────────────────────────────────────────────────────────────
\title{Digital Twin Enabled Sustainability Intelligence Framework
for Smart Cities: A Case Study of Delhi%
\thanks{This work was conducted as part of the InfraVision Project,
Urban AI and Sustainability Systems research initiative, B.Tech
Computer Science and Engineering Programme, 2026.}}

\author{%
\IEEEauthorblockN{Asad Ahmad}
\IEEEauthorblockA{%
\textit{Department of Computer Science and Engineering}\\
\textit{Urban AI \& Sustainability Systems --- InfraVision}\\
New Delhi, India\\
techieasad01@gmail.com}
\and
\IEEEauthorblockN{Kislay Kumar}
\IEEEauthorblockA{%
\textit{Department of Computer Science and Engineering}\\
\textit{Urban AI \& Sustainability Systems --- InfraVision}\\
New Delhi, India\\
kislay@example.com}
}

\maketitle

% ─────────────────────────────────────────────────────────────────────────────
%  ABSTRACT
% ─────────────────────────────────────────────────────────────────────────────
\begin{abstract}
Rapid urbanization in Indian megacities has produced compounding and
interdependent stresses across water supply, energy distribution, solid
waste management, and atmospheric emissions simultaneously. Conventional
planning instruments address these domains in isolation, precluding the
cross-domain reasoning that sustainable urban governance requires. This
paper presents InfraVision, a Digital Twin enabled sustainability
intelligence framework applied to the National Capital Territory of
Delhi. The system represents the city as a nine-zone directed multi-edge
graph and integrates four domain-specific machine learning forecasting
models --- Prophet, XGBoost, Random Forest, and Ridge Regression ---
with a PuLP-based constraint-driven binary linear optimisation engine
and a SHAP-driven explainability layer. Baseline projections to 2030
indicate that the water supply deficit will widen from 594 to 764
million gallons per day, greenhouse gas emissions will increase from
56.5 to 62.4 MtCO2e, and waste processing efficiency will improve from
78 to 95 percent. Under an optimised policy bundle satisfying fiscal and
emission constraints, the composite sustainability score improves by
5.08 points and carbon emissions decline by 9.49 MtCO2e relative to
baseline. A four-component validation protocol yields mean absolute
percentage errors below six percent across all forecasting domains,
confirms directional consistency of simulation outputs with the
international evidence base, demonstrates input robustness through
sensitivity analysis, and shows that InfraVision reduces decision cycle
time by more than 95 percent while delivering a 3.6-point additional
sustainability improvement over conventional planning workflows under
equivalent budget assumptions.
\end{abstract}

\begin{IEEEkeywords}
digital twin, urban sustainability, smart city, machine learning
forecasting, policy optimisation, explainable AI, Delhi, carbon
intelligence
\end{IEEEkeywords}

% ─────────────────────────────────────────────────────────────────────────────
%  I. INTRODUCTION
% ─────────────────────────────────────────────────────────────────────────────
\section{Introduction}

Cities are responsible for a disproportionate share of global resource
consumption and greenhouse gas emissions, and this share is growing.
Urban populations are projected to represent nearly 70 percent of the
global total by 2050 \cite{un2023}, intensifying pressures on water
supply systems, energy grids, waste processing facilities, and
atmospheric quality simultaneously. Delhi, with a resident population
exceeding 33 million and a persistently high in-migration rate,
exemplifies this compounding condition. The Delhi Jal Board reported a
daily water deficit in excess of 300 million gallons per day (MGD) as
early as 2023 \cite{djb2023}; per capita greenhouse gas emissions
continue to rise as the vehicle fleet and industrial base expand; and
waste processing capacity at the city's landfill sites is approaching
functional saturation \cite{cpcb2023}.

Conventional urban planning relies on periodic, domain-siloed
assessments that cannot capture the interdependencies between these
systems or support anticipatory policy design. Traditional approaches
are generally too inflexible and reactive to manage the complexity and
dynamism of modern cities \cite{therias2023}. Digital twin technology
offers a structural answer to this problem by creating a continuously
updated virtual replica of the physical city that supports real-time
monitoring, predictive analytics, and policy simulation
\cite{mazzetto2024}. When applied to urban sustainability, a digital
twin can model how a policy intervention in one domain --- for example,
expanding waste-to-energy capacity --- propagates its effects across
energy supply, emission accounting, and spatial equity simultaneously.

Recent advances in the Internet of Things (IoT), artificial
intelligence, geographic information systems, and cloud computing have
accelerated digital twin deployment across smart city development
\cite{alexandrov2024}. However, the urban sustainability literature
reveals three persistent gaps: most implementations are domain-specific
and do not model cross-domain feedback; optimisation is rarely embedded
in the twin itself; and the explainability of model output is largely
absent, limiting the trust that planners can place in data-driven
recommendations \cite{alkhateeb2024}.

This paper presents InfraVision, an integrated Digital Twin enabled
sustainability intelligence framework that addresses each of these gaps.
The system makes five specific contributions to the literature: (i) an
integrated multi-domain digital twin jointly modelling water, energy,
waste, and carbon; (ii) a constraint-based linear optimisation engine
for auditable policy ranking; (iii) a SHapley Additive exPlanations
(SHAP) \cite{lundberg2017}-driven explainability layer;
(iv) an independent four-component
validation protocol; and (v) a detailed empirical analysis of Delhi's
sustainability trajectory to 2030.

The remainder of this paper is organised as follows.
Section~\ref{sec:related} reviews the related literature.
Section~\ref{sec:novelty} articulates the novelty and specific
contributions of the proposed system relative to existing work.
Section~\ref{sec:method} presents the methodology.
Section~\ref{sec:arch} describes the system architecture.
Section~\ref{sec:results} presents the empirical results.
Section~\ref{sec:validation} provides the validation and performance
evaluation. Section~\ref{sec:discussion} discusses the implications.
Section~\ref{sec:conclusion} concludes.

% ─────────────────────────────────────────────────────────────────────────────
%  II. RELATED WORK
% ─────────────────────────────────────────────────────────────────────────────
\section{Related Work}
\label{sec:related}

\subsection{Urban Digital Twins}

The digital twin concept originated in manufacturing, where it was used
to monitor performance and predict failure of physical assets. Translated
to the urban scale, an urban digital twin (UDT) is a dynamic,
continuously updated virtual representation of a city's physical systems,
integrating data from sensors, IoT devices, geographic information
systems, and administrative databases \cite{alkhateeb2024}. Unlike
static spatial models, UDTs enable real-time monitoring, scenario
simulation, and adaptive decision-making \cite{therias2023}. Cities are
treated as complex adaptive systems in which physical, digital, and
social subsystems interact nonlinearly, and the twin provides a cyber
representation of these interactions to support predictive governance
\cite{alexandrov2024}.

\subsection{Digital Twins for Sustainability}

Sustainability research has increasingly recognised the digital twin as
a practical instrument for managing urban resource flows.
Al-Khateeb et al. \cite{alkhateeb2024} demonstrate that digital twins
integrating energy, water, and waste data can detect inefficiencies and
support circular economy interventions. Mazzetto \cite{mazzetto2024}
shows that AI-driven twins can reduce building energy consumption and
carbon emissions through real-time adaptive control. Climate-focused
implementations have addressed flood risk modelling, urban heat island
mitigation, and air quality management \cite{zhu2025}. Several studies
link UDT outputs to Sustainable Development Goal monitoring, particularly
SDG 11 (Sustainable Cities) and SDG 13 (Climate Action).

\subsection{Machine Learning in Urban Forecasting}

Prophet \cite{taylor2018} is well-suited to municipal demand series with
strong seasonal patterns and trend changepoints. XGBoost \cite{chen2016}
handles non-linear interactions among temperature, occupancy, and tariff
variables characteristic of urban energy demand. Random Forest provides
robustness to heterogeneous and partially missing data typical of
ward-level waste records. Ridge Regression offers stable regularised
estimates for emission inventories where multicollinearity among activity
drivers is common.

\subsection{Research Gap}

A systematic review of the UDT sustainability literature identifies three
persistent gaps: (i) domain-siloed modelling that ignores cross-domain
feedback; (ii) heuristic rather than optimisation-driven policy
selection; and (iii) absence of explainability mechanisms that would
enable planners to interrogate model recommendations. InfraVision
addresses each gap explicitly.

% ─────────────────────────────────────────────────────────────────────────────
%  III. NOVELTY AND KEY CONTRIBUTIONS
% ─────────────────────────────────────────────────────────────────────────────
\section{Novelty and Key Contributions}
\label{sec:novelty}

The urban digital twin literature, while extensive, has not produced a
system that simultaneously addresses multi-domain integration,
optimisation-driven policy selection, quantified explainability, and
independent statistical validation within a single operational
framework. InfraVision is specifically designed to fill this gap.
Table~\ref{tab:novelty} positions the system against five representative
prior works across the four defining dimensions.

\begin{table}[htbp]
\caption{Novelty Positioning: InfraVision vs. Prior Digital Twin Systems}
\label{tab:novelty}
\begin{center}
\begin{tabular}{|p{1.6cm}|c|c|c|c|}
\hline
\textbf{System} &
\textbf{\shortstack{Multi-\\Domain}} &
\textbf{\shortstack{LP\\Optim.}} &
\textbf{\shortstack{Explain-\\ability}} &
\textbf{\shortstack{Formal\\Validation}} \\
\hline
Therias \& Rafiee \cite{therias2023}   & No  & No  & No  & No  \\
Mazzetto \cite{mazzetto2024}           & No  & No  & No  & No  \\
Alexandrov et al. \cite{alexandrov2024}& Partial & No & No & Partial \\
Al-Khateeb et al. \cite{alkhateeb2024} & Partial & No & No & No  \\
Zhu \& Jin \cite{zhu2025}              & No  & No  & No  & No  \\
\textbf{InfraVision (Ours)}            & \textbf{Yes} & \textbf{Yes} &
\textbf{Yes} & \textbf{Yes} \\
\hline
\end{tabular}
\end{center}
\end{table}

Five specific contributions differentiate InfraVision from prior art:

\textbf{C1 --- Integrated Four-Domain Digital Twin.}
Existing UDT implementations address at most two sustainability domains
concurrently. InfraVision jointly models water supply, energy demand,
solid waste, and carbon emissions on a shared directed graph, enabling
the detection of cross-domain trade-offs that domain-siloed systems
cannot observe. The cascade failure analysis of Section~\ref{sec:results}
is a direct product of this integration.

\textbf{C2 --- Constraint-Based Policy Optimisation.}
Policy selection in prior UDT systems is invariably heuristic ---
relying on expert judgement or manually ranked priority lists.
InfraVision embeds a binary linear programme that simultaneously
satisfies a capital budget constraint, a minimum emission reduction
target, and a minimum sustainability score threshold, producing an
auditable, reproducible policy ranking rather than an opinion.

\textbf{C3 --- Quantified Explainability via SHAP.}
AI-driven urban planning recommendations are frequently opaque,
eroding planner trust. InfraVision applies SHAP feature attribution to
decompose every composite score prediction into per-driver contributions,
providing a transparent account of why the model recommends a particular
intervention and which input assumptions are most consequential.

\textbf{C4 --- Formal Four-Component Validation Protocol.}
Most UDT papers report only in-sample accuracy or qualitative
plausibility arguments. InfraVision submits every forecasting component
to a held-out empirical test, validates simulation direction against
documented international case studies, characterises robustness through
systematic input perturbation, and benchmarks decision quality against a
representative traditional workflow --- a protocol that to the best of
the authors' knowledge has not previously been applied to a
sustainability-focused urban digital twin.

\textbf{C5 --- First Multi-Domain Sustainability Digital Twin for Delhi.}
Despite Delhi being among the world's most stressed urban environments
along multiple sustainability dimensions simultaneously, no prior study
has modelled the joint water--energy--waste--carbon system of the
National Capital Territory within a unified digital twin. InfraVision
provides this foundation and makes its empirical findings available as
a reproducible baseline for future comparative studies.

% ─────────────────────────────────────────────────────────────────────────────
%  IV. METHODOLOGY
% ─────────────────────────────────────────────────────────────────────────────
\section{Methodology}
\label{sec:method}

\subsection{Digital Twin Graph Representation}

The city is modelled as a directed multi-edge graph
$\mathcal{G} = (\mathcal{V}, \mathcal{E})$, where $\mathcal{V}$ denotes
the nine administrative zones of Delhi and $\mathcal{E}$ denotes
infrastructure linkages connecting them. Formally:

\begin{equation}
\mathcal{G} = \bigl(\{v_1, \ldots, v_9\},\; \mathcal{E}^w \cup
\mathcal{E}^e \cup \mathcal{E}^{wst} \cup \mathcal{E}^c\bigr)
\label{eq:graph}
\end{equation}

where $\mathcal{E}^w$, $\mathcal{E}^e$, $\mathcal{E}^{wst}$, and
$\mathcal{E}^c$ denote edge sets typed by the water, energy, waste, and
carbon domains respectively. Each edge carries attributes for
transmission capacity, loss coefficient, and service latency.
Zone criticality $\kappa_i$ is defined as the weighted sum of downstream
impact across all edge types:

\begin{equation}
\kappa_i = \sum_{j \in \mathcal{N}(i)} \sum_{d \in D} w_d \cdot
\ell_{ij}^d
\label{eq:criticality}
\end{equation}

where $\mathcal{N}(i)$ is the set of zones reachable from zone $i$,
$D = \{w, e, wst, c\}$ is the domain set, $w_d$ is a domain importance
weight, and $\ell_{ij}^d$ is the load transferred along edge
$(i,j)$ in domain $d$.

\subsection{Machine Learning Forecasting Layer}

Four domain-specific models were selected on the basis of the
statistical character of their target series:

\begin{itemize}
\item \textbf{Prophet} \cite{taylor2018}: Water demand forecasting.
Captures trend changepoints, weekly and annual seasonality, and
scheduled intervention effects (e.g., reservoir maintenance).
\item \textbf{XGBoost} \cite{chen2016}: Energy demand forecasting.
Handles non-linear interactions among temperature, day-of-week,
occupancy, and retail tariff variables without manual feature
engineering.
\item \textbf{Random Forest}: Solid waste generation forecasting.
Tolerates heterogeneous, partially missing ward-level inputs and
provides ensemble-based variance estimates.
\item \textbf{Ridge Regression}: Carbon emission inventory modelling.
Controls multicollinearity among correlated activity drivers through
$L_2$ regularisation.
\end{itemize}

All models were trained on 2015--2022 data and evaluated on a withheld
18-month test set (2023--mid 2024).

\subsection{Constraint-Based Optimisation Model}

Policy selection is cast as a binary linear programme solved via the
CBC solver through PuLP. Let $I_i \in \{0,1\}$ denote the intervention
decision for policy $i$, $E_i$ its expected composite sustainability
impact, $c_i$ its capital cost, and $r_i$ its per-unit greenhouse gas
reduction. The programme is:

\begin{equation}
\text{Maximise}\quad Z = \sum_{i=1}^{n} E_i \cdot I_i
\label{eq:obj}
\end{equation}

subject to:

\begin{align}
\sum_{i=1}^{n} c_i I_i &\leq B \label{eq:budget}\\
\sum_{i=1}^{n} r_i I_i &\geq R^* \label{eq:emission}\\
\Delta S &\geq S^* \label{eq:score}\\
I_i &\in \{0,1\},\quad \forall\, i \label{eq:binary}
\end{align}

where $B$ is the available capital budget, $R^*$ is the minimum required
emission reduction, $S^*$ is the minimum sustainability score
improvement, and $\Delta S$ is the score change delivered by the
selected bundle. This formulation ensures recommendations are
simultaneously fiscally feasible and emission-compliant.

\subsection{Data Sources and Pre-processing}

Historical inputs were assembled from four primary sources: the Delhi
Jal Board for zone-level water demand and supply records (2015--2023)
\cite{djb2023}; the Central Pollution Control Board for annual greenhouse
gas inventories (2015--2023) \cite{cpcb2023}; the Municipal Corporation
of Delhi for ward-level solid waste records (2015--2023); and the Delhi
State Load Despatch Centre for hourly electricity demand profiles
(2018--2023). Missing values were imputed via seasonal decomposition for
gaps exceeding two weeks and by zone-level mean for shorter gaps. All
features were standardised prior to model training, and graph attributes
were calibrated against the 2023 infrastructure audit of the Delhi Urban
Development Authority.

% ─────────────────────────────────────────────────────────────────────────────
%  IV. SYSTEM ARCHITECTURE
% ─────────────────────────────────────────────────────────────────────────────
\section{System Architecture and Implementation}
\label{sec:arch}

InfraVision follows a four-layer modular architecture that separates
data acquisition, analytical modelling, intelligent explanation, and
interactive presentation. Table~\ref{tab:arch} summarises the layers,
technology stack, and primary function of each.

\begin{table}[htbp]
\caption{InfraVision System Architecture}
\label{tab:arch}
\begin{center}
\begin{tabular}{|p{1.3cm}|p{2.4cm}|p{3.8cm}|}
\hline
\textbf{Layer} & \textbf{Technology} & \textbf{Function} \\
\hline
Data Ingestion & Python, SQLAlchemy & Acquisition, cleaning, feature engineering \\
\hline
Analytics & NetworkX, Prophet, XGBoost, sklearn, PuLP & Graph modelling, domain forecasting, optimisation \\
\hline
Intelligence & SHAP, Python & Feature attribution and explainability \\
\hline
Presentation & Next.js, TypeScript & Interactive planner-facing decision support \\
\hline
\end{tabular}
\end{center}
\end{table}

The end-to-end data flow proceeds through five sequential stages: (i)
raw data acquisition and schema validation; (ii) feature engineering and
graph attribute calibration; (iii) domain forecasting and scenario
projection; (iv) constraint-based policy optimisation; and (v) SHAP
attribution and dashboard presentation. The modular design ensures that
any individual component --- for example, the energy forecasting model
--- can be retrained or replaced without disrupting the remainder of the
pipeline.

% ─────────────────────────────────────────────────────────────────────────────
%  V. RESULTS AND ANALYSIS
% ─────────────────────────────────────────────────────────────────────────────
\section{Results and Analysis}
\label{sec:results}

\subsection{Baseline Sustainability Projections (2025--2030)}

The forecasting layer produces five-year business-as-usual projections
across all four sustainability domains. Table~\ref{tab:baseline}
presents the headline indicators.

\begin{table}[htbp]
\caption{Baseline Sustainability Projections, 2025--2030}
\label{tab:baseline}
\begin{center}
\begin{tabular}{|l|c|c|c|}
\hline
\textbf{Indicator} & \textbf{2025} & \textbf{2030} & \textbf{Trend} \\
\hline
Water Deficit (MGD)       & 594   & 764   & \textbf{Worsening} \\
Waste Processing Eff. (\%)& 78    & 95    & Improving \\
GHG Emissions (MtCO$_2$e) & 56.5  & 62.4  & \textbf{Worsening} \\
Sustainability Score      & 60.22 & 62.44 & Improving \\
\hline
\end{tabular}
\end{center}
\end{table}

The water gap is projected to widen by 170 MGD, driven by population
growth in peripheral zones and a distribution network leakage rate
estimated at 36 percent of total supply. Greenhouse gas emissions are
expected to grow at 2.0 percent per annum on average, with the transport
and energy sectors contributing approximately 72 percent of the total
inventory. Waste processing efficiency improves most markedly in the
central and southern zones, where new processing capacity has been
committed, while peripheral zones remain significantly below the city
average. The modest aggregate sustainability score gain masks a
structurally divided trajectory: two domains improve while two
deteriorate.

\subsection{Policy Simulation Outcomes}

The optimisation engine (equations \eqref{eq:obj}--\eqref{eq:binary})
evaluated 14 candidate policy interventions against the stated
constraints. The optimal bundle consists of four instruments:
non-revenue water reduction investments, waste-to-energy plant
expansion, a 40 percent electric vehicle fleet transition, and a
rooftop solar incentive scheme. This bundle yields:

\begin{itemize}
\item Composite sustainability score improvement: \textbf{+5.08 points}
above the 2030 baseline.
\item Carbon emission reduction: \textbf{$-$9.49 MtCO$_2$e} relative
to baseline.
\end{itemize}

These results confirm that structured, constraint-based policy design
substantially outperforms the business-as-usual trajectory within
realistic fiscal envelopes.

\subsection{Critical Infrastructure Analysis}

Zone criticality scoring via equation \eqref{eq:criticality} reveals
that the Central zone carries the highest single-node risk in the graph.
Its cascading failure propagates across four directly dependent zones
within the water domain, representing 47 percent of the city's total
residential supply network. Maintenance and redundancy investments in
the Central zone therefore produce disproportionately large system-level
resilience gains relative to equivalent investments elsewhere in the
network, a finding invisible to domain-siloed analysis.

% ─────────────────────────────────────────────────────────────────────────────
%  VI. SYSTEM VALIDATION AND PERFORMANCE EVALUATION
% ─────────────────────────────────────────────────────────────────────────────
\section{System Validation and Performance Evaluation}
\label{sec:validation}

A predictive planning system is only as useful as the confidence that
can be placed in its outputs. This section presents a structured,
four-component validation protocol applied to the InfraVision digital
twin, following principles recommended for UDT validation
\cite{alexandrov2024} and for assessing the decision usefulness of
AI-driven urban systems \cite{mazzetto2024}.

\subsection{Validation Strategy}

Validation was designed around the principle that a digital twin must
reproduce observable sustainability indicators before its forward
projections can be trusted for policy purposes. Three indicator families
were selected: water demand, greenhouse gas emissions, and waste
processing throughput. The protocol proceeds through four stages.
First, historical ground truth was assembled from the Delhi Jal Board,
Central Pollution Control Board, and Municipal Corporation of Delhi for
the period 2015--2024. Second, the twin was operated in hindcast mode
over the 2015--2022 training window, and predictions for 2023--2024 were
compared with withheld data. Third, residual errors were decomposed by
zone and season to detect systematic spatial or temporal bias. Fourth,
predicted trajectories were overlaid against observed series to verify
that the model reproduces both the level and the trend shape of each
indicator --- not only isolated point values.

\subsection{Model Evaluation Metrics}
\label{subsec:metrics}

Forecasting accuracy was quantified using Root Mean Squared Error
(RMSE), Mean Absolute Error (MAE), and Mean Absolute Percentage Error
(MAPE). All metrics were computed on the withheld 18-month test set.
Results are presented in Table~\ref{tab:metrics}.

\begin{table}[htbp]
\caption{Forecasting Model Evaluation Metrics (18-Month Test Set)}
\label{tab:metrics}
\begin{center}
\begin{tabular}{|l|l|c|c|c|}
\hline
\textbf{Domain} & \textbf{Model} & \textbf{RMSE} & \textbf{MAE} &
\textbf{MAPE} \\
\hline
Water Demand   & Prophet         & 12.4 MGD          & 9.1 MGD         & 3.8\% \\
Energy Demand  & XGBoost         & 84.7 MW            & 61.3 MW         & 4.2\% \\
Waste Gen.     & Random Forest   & 112 t/day          & 87 t/day        & 5.6\% \\
Carbon Emiss.  & Ridge Reg.      & 0.41 MtCO$_2$e    & 0.32 MtCO$_2$e & 5.1\% \\
\hline
\end{tabular}
\end{center}
\end{table}

All four models report MAPE values below six percent, which is within
the tolerance conventionally accepted for municipal demand planning at
five-year horizons. Prophet achieves the lowest MAPE of 3.8 percent
because water demand in Delhi is strongly seasonal and trend-dominated,
aligning with the model's inductive bias. XGBoost delivers 4.2 percent
MAPE on energy demand by capturing non-linear temperature and occupancy
interactions. The Random Forest model tolerates the heterogeneous
ward-level waste inputs and achieves 5.6 percent MAPE without
overfitting. Ridge Regression provides a stable regularised estimate for
carbon where multicollinearity among activity drivers is most severe.
Residual analysis confirms the absence of systematic seasonal or spatial
bias in all four models.

\subsection{Simulation Validation}

The simulation layer must produce policy outcomes that are directionally
consistent with established urban planning theory and with documented
responses in comparable cities. Four scenarios were simulated and
compared with expected trends derived from peer-reviewed case studies
of Singapore, Helsinki, Seoul, and Amsterdam. Table~\ref{tab:simval}
summarises the assessment.

\begin{table}[htbp]
\caption{Simulation Validation: Outcomes vs. Expected Real-World Trends}
\label{tab:simval}
\begin{center}
\begin{tabular}{|p{2.0cm}|p{2.2cm}|p{2.0cm}|c|}
\hline
\textbf{Policy} & \textbf{Simulated} & \textbf{Expected} &
\textbf{Result} \\
\hline
Waste-to-Energy Expansion &
  $+$14\% processing, $-$2.1 MtCO$_2$e &
  Processing uplift, moderate emission reduction &
  Pass \\
\hline
Non-Revenue Water Reduction &
  $-$118 MGD deficit &
  Substantial deficit reduction within 5 yr &
  Pass \\
\hline
EV Fleet (40\%) &
  $-$3.4 MtCO$_2$e, $+$6\% grid load &
  Emission fall, grid stress &
  Pass \\
\hline
Rooftop Solar &
  $-$1.8 MtCO$_2$e, $+$2.1 pts &
  Gradual emission decline, score uplift &
  Pass \\
\hline
\end{tabular}
\end{center}
\end{table}

Every simulated policy moves target indicators in the direction
predicted by the international evidence base. No scenario produces
paradoxical results --- such as emission reduction without energy
substitution, or water gap closure without loss reduction. The
simulation layer therefore satisfies the minimum criterion of
directional plausibility in addition to the numerical accuracy
demonstrated in Section~\ref{subsec:metrics}.

\subsection{Sensitivity Analysis}
\label{subsec:sensitivity}

Sensitivity analysis quantifies how strongly system outputs respond to
perturbation of uncertain inputs. Three drivers were perturbed
independently across a $\pm$10 percent range in 2-percent increments,
with all other parameters held fixed: projected population growth rate,
per capita water demand, and per capita energy demand. Output elasticity
$\varepsilon$ is defined as:

\begin{equation}
\varepsilon = \frac{\partial \hat{y}/\hat{y}}{\partial x/x}
\label{eq:elasticity}
\end{equation}

Results are summarised in Table~\ref{tab:sensitivity}.

\begin{table}[htbp]
\caption{Sensitivity Analysis: Output Elasticities}
\label{tab:sensitivity}
\begin{center}
\begin{tabular}{|l|c|c|}
\hline
\textbf{Perturbed Input} & \textbf{Primary Output} &
\textbf{Elasticity $\varepsilon$} \\
\hline
Population Projection     & Water deficit     & 0.96 \\
Per Capita Water Demand   & Water deficit     & 1.03 \\
Per Capita Energy Demand  & GHG emissions     & 0.74 \\
Waste Generation Rate     & Processing rate   & 0.24 \\
\hline
\end{tabular}
\end{center}
\end{table}

Elasticities for water-related drivers are close to unity, indicating
that projections scale proportionally with demographic uncertainty --- a
behaviour consistent with a well-calibrated physical demand model.
Energy and waste domains exhibit lower elasticities, reflecting
stabilising effects of installed capacity constraints and tariff
inertia. Crucially, no perturbation within the tested range produces a
sign reversal in any headline indicator, establishing that the framework
is robust: realistic errors in demographic or demand inputs propagate
proportionally and do not generate qualitatively misleading planning
recommendations.

\subsection{Benchmark Comparison}

To contextualise the improvement offered by InfraVision, its decision
quality was compared against a traditional planning workflow
representative of current Indian municipal planning practice --- one
based on annual static reports, domain-siloed spreadsheets, and
heuristic policy selection. The comparison is presented in
Table~\ref{tab:benchmark}.

\begin{table}[htbp]
\caption{Benchmark Comparison: InfraVision vs. Traditional Planning}
\label{tab:benchmark}
\begin{center}
\begin{tabular}{|p{2.0cm}|p{2.0cm}|p{2.2cm}|}
\hline
\textbf{Dimension} & \textbf{Traditional} & \textbf{InfraVision} \\
\hline
Decision cycle time & 6--8 weeks & Under 2 hours \\
Domains considered & 1--2 & 4 (water, energy, waste, carbon) \\
Policy ranking basis & Heuristic & Constraint-based LP \\
Output explainability & Qualitative & SHAP attributions \\
Cascade risk detection & Not performed & Graph criticality score \\
Sustainability score uplift & $+$1.4 pts (observed) & $+$5.08 pts (optimised) \\
\hline
\end{tabular}
\end{center}
\end{table}

The gains extend well beyond faster computation. InfraVision enables
joint reasoning across four sustainability domains, produces auditable
constraint-based rankings, introduces cascade risk detection as a
capability absent from the traditional workflow entirely, and delivers a
3.6-point additional improvement in composite sustainability score under
identical budget assumptions. These results are consistent with
Zhu and Jin \cite{zhu2025}, who find that AI-augmented UDTs
consistently outperform conventional planning instruments on composite
resilience and sustainability metrics.

\subsection{Validation Summary}

The four-component protocol produces convergent evidence. The
forecasting layer achieves domain-level MAPE values of 3.8--5.6 percent
on held-out data. The simulation layer is directionally consistent with
the international evidence base across all tested scenarios. The
framework is robust to realistic input perturbations, with elasticities
that remain proportional and sign-stable. Against a traditional
planning workflow, InfraVision delivers superior coverage, auditability,
and composite sustainability performance. These results confirm that the
system is not only technically sound but decision-useful: it produces
intelligence that municipal planners can act upon with quantified
confidence.

% ─────────────────────────────────────────────────────────────────────────────
%  VII. DISCUSSION
% ─────────────────────────────────────────────────────────────────────────────
\section{Discussion}
\label{sec:discussion}

\subsection{Interpretation of Key Findings}

Four substantive observations emerge from the combined empirical and
validation results. First, the water supply deficit is structurally the
most critical sustainability challenge in the Delhi system. The projected
170 MGD widening between 2025 and 2030 implies that a supply-side or
loss-reduction intervention of equivalent scale must be commissioned
merely to prevent the gap from growing further --- before any absolute
improvement is achievable. The optimisation layer identifies non-revenue
water reduction as the single highest-return intervention in the
14-candidate set, yielding 118 MGD of effective supply gain per unit of
capital committed. This finding aligns with the broader evidence that
distribution loss reduction offers superior returns relative to new
source development in mature urban networks \cite{alkhateeb2024} and
underscores that infrastructure efficiency, not only capacity expansion,
must anchor Delhi's water strategy.

Second, the divergence between improving waste performance and worsening
greenhouse gas emissions exposes a structural limitation of aggregate
sustainability scoring. The composite score registers a modest 2.22-point
gain under business-as-usual conditions because waste processing gains
numerically outweigh water and emission losses in the current weighting
scheme. This arithmetic improvement conceals the fact that the two most
ecologically consequential sub-indices are simultaneously deteriorating.
Any policy framework that optimises solely for the aggregate score
without imposing hard sub-index constraints will generate misleading
performance narratives and politically convenient but environmentally
inadequate outcomes. The binding emission constraint in
equation~\eqref{eq:emission} is the mechanism that prevents this failure
mode within InfraVision.

Third, the cascade risk analysis reveals a systemic vulnerability in the
Central zone that no domain-level indicator surfaced independently. The
zone's failure propagates across 47 percent of the city's residential
water supply network, and this cross-zone coupling is a property of the
infrastructure topology, not of any individual domain's performance
metric. This finding exemplifies Contribution C1: multi-domain joint
modelling generates qualitatively different --- and more actionable ---
planning intelligence than domain-siloed analysis.

Fourth, the validation results establish an important asymmetry: the
system is most accurate where the physical dynamics are most regular
(water demand, MAPE 3.8\%) and least accurate where they are most
heterogeneous (waste generation, MAPE 5.6\%). This ordering is
theoretically predictable and confirms that model selection was
appropriate; a single universal model would have performed worse across
the full range.

\subsection{Limitations}

Three limitations bound the current study. The graph model reflects a
2023 infrastructure audit snapshot and does not yet incorporate
real-time IoT telemetry; drift in network topology or capacity since
that audit will not be reflected in current projections. The binary
linear optimisation assumes that policy impacts are additive and
independent, which may overestimate the combined benefit of the optimal
bundle when simultaneous implementation competes for the same labour,
supply chains, or political capital. Finally, the generalisability of
both the model parameterisation and the policy rankings to other Indian
megacities has not been empirically verified; Delhi-specific
infrastructure configurations, leakage rates, and emissions profiles may
not transfer without re-calibration.

% ─────────────────────────────────────────────────────────────────────────────
%  VIII. CONCLUSION
% ─────────────────────────────────────────────────────────────────────────────
\section{Conclusion}
\label{sec:conclusion}

This paper has presented InfraVision, a Digital Twin enabled
sustainability intelligence framework for smart city planning, and
evaluated it on the case of Delhi. Five specific contributions
distinguish the system from prior work: integrated four-domain
modelling; constraint-based binary linear policy optimisation; SHAP
explainability; a formal four-component validation protocol; and the
first multi-domain sustainability digital twin for the National Capital
Territory.

Empirically, baseline projections reveal a structurally divided
trajectory --- waste processing efficiency improves to 95 percent while
the water deficit widens to 764 MGD and greenhouse gas emissions rise to
62.4 MtCO$_2$e by 2030. Under the optimised policy bundle, the composite
sustainability score improves by 5.08 points and carbon emissions decline
by 9.49 MtCO$_2$e, demonstrating that constraint-based planning
substantially outperforms the business-as-usual trajectory within
realistic fiscal envelopes. Cascade risk analysis identifies the Central
zone as the highest single-node vulnerability in the infrastructure
network --- a structural finding invisible to any domain-siloed
approach.

The four-component validation protocol yields MAPE values of 3.8--5.6
percent across all forecasting domains on withheld test data, confirms
directional consistency of all simulated policy outcomes with the
international evidence base, establishes sign-stable input elasticities,
and demonstrates a decision cycle reduction of more than 95 percent with
a 3.6-point additional sustainability improvement over conventional
planning workflows.

Future work will pursue three extensions: integration of live IoT
telemetry to move the twin from snapshot-based to continuous operation;
generalisation to additional Indian megacities (Mumbai, Bengaluru,
Hyderabad) to validate transferability; and augmentation of the
optimisation layer with reinforcement learning to capture diminishing
returns and dynamic interaction effects when multiple interventions are
deployed simultaneously.

% ─────────────────────────────────────────────────────────────────────────────
%  ACKNOWLEDGMENT
% ─────────────────────────────────────────────────────────────────────────────
\section*{Acknowledgment}

The authors acknowledge the publicly available sustainability data
provided by the Delhi Jal Board, the Central Pollution Control Board,
and the Municipal Corporation of Delhi, without which the empirical
analysis in this paper would not have been possible. The authors also
thank the open-source communities behind NetworkX, PuLP, Prophet,
XGBoost, scikit-learn, and SHAP for maintaining the libraries on which
this system is built.

% ─────────────────────────────────────────────────────────────────────────────
%  REFERENCES
% ─────────────────────────────────────────────────────────────────────────────
\begin{thebibliography}{00}

\bibitem{un2023}
United Nations, ``World Urbanization Prospects: The 2023 Revision,''
Department of Economic and Social Affairs, New York, USA, 2023.

\bibitem{djb2023}
Delhi Jal Board, ``Annual Report on Water Supply and Distribution
2022--23,'' Government of the National Capital Territory of Delhi,
New Delhi, India, 2023.

\bibitem{cpcb2023}
Central Pollution Control Board, ``National Emission Inventory for
Indian Megacities 2022--23,'' Ministry of Environment, Forest and
Climate Change, Government of India, New Delhi, 2023.

\bibitem{therias2023}
E. Therias and A. Rafiee, ``Urban digital twins in practice: A critical
review of planning applications,'' \textit{Cities}, vol. 134,
p. 104168, 2023.

\bibitem{mazzetto2024}
S. Mazzetto, ``Digital twins for smart urban systems: Architecture,
applications, and challenges,'' \textit{Smart Cities}, vol. 7, no. 2,
pp. 445--470, 2024.

\bibitem{alexandrov2024}
D. Alexandrov, J. Kim, and I. Petrov, ``Enabling technologies for urban
digital twins: IoT, AI, and geospatial integration,''
\textit{Smart Cities}, vol. 7, no. 1, pp. 112--134, 2024.

\bibitem{alkhateeb2024}
A. Al-Khateeb, M. Yessef, and A. Hassan, ``Urban digital twins for
sustainable resource management: A review,''
\textit{Sustainable Cities and Society}, vol. 98, p. 104801, 2024.

\bibitem{zhu2025}
S. Zhu and Y. Jin, ``Climate resilience and urban digital twins: A
systematic review,'' \textit{Environmental Modelling and Software},
vol. 178, p. 105801, 2025.

\bibitem{taylor2018}
S. J. Taylor and B. Letham, ``Forecasting at scale,''
\textit{The American Statistician}, vol. 72, no. 1, pp. 37--45, 2018.

\bibitem{chen2016}
T. Chen and C. Guestrin, ``XGBoost: A scalable tree boosting system,''
in \textit{Proc. 22nd ACM SIGKDD Int. Conf. Knowledge Discovery and
Data Mining}, San Francisco, CA, USA, 2016, pp. 785--794.

\bibitem{yessef2025}
M. Yessef \textit{et al.}, ``Digital twin technology in smart cities:
Architecture, applications, and research directions,''
\textit{Int. J. Urban Systems}, vol. 3, no. 1, pp. 22--49, 2025.

\bibitem{lim2025}
F. P. C. Lim, ``Smart city planning through digital twins: Enhancing
urban resilience and sustainability,''
\textit{Int. J. Science, Architecture, Technology and Environment},
vol. 2, no. 11, pp. 140--152, Nov. 2025.

\bibitem{lundberg2017}
S. M. Lundberg and S.-I. Lee, ``A unified approach to interpreting
model predictions,'' in \textit{Proc. Advances in Neural Information
Processing Systems (NeurIPS)}, Long Beach, CA, USA, 2017, pp.
4765--4774.

\bibitem{mcd2023}
Municipal Corporation of Delhi, ``Solid Waste Management Annual Review
2022--23,'' MCD Press, New Delhi, India, 2023.

\end{thebibliography}

\end{document}
